% GENERATED FILE. Edit canonical lessons, then run npm run generate:books.
% canonical-source-hash: fnv1a64:33940d9d365bdfd2
% canonical-lessons: MR-C510-raja-leave, MR-C510-bekari, MR-C510-rajinama-dene, MR-C510-nivrtta-hone, MR-C510-badtarph-karne

\chapter{Leaving Work}
\label{ch:mr-leaving-work}
\hlchaptermodality{510}

\begin{chapteropening}
\textbf{By the end of this chapter:} \emph{I can talk about leave, strikes, resigning and retiring.}

Complete the last lesson of chapter 510: I can talk about leave, strikes, resigning and retiring.
\end{chapteropening}

\section[\texorpdfstring{\mr{रजा} (rajā)}{रजा (rajā)}]{\mr{रजा} (rajā) — time off work; leave}

\label{lesson:MR-C510-raja-leave}



\begin{quote}
Before the new one: say the Marathi for junior, then the Marathi for to strive.
\end{quote}

\subsection*{You'll want to know: \mr{रजा}}
\textbf{\mr{रजा}} — \emph{rajā} — \textquotedblleft{}time off work; leave\textquotedblright{}.

\begin{quote}
\mr{आजारी} \mr{असल्यामुळे} \mr{मी} \mr{दोन} \mr{दिवस} \mr{रजा} \mr{घेतली}.
\end{quote}

\begin{quote}
\emph{ājārī aslyāmuḷe mī don divas rajā ghetlī.}
\end{quote}

\begin{quote}
Because I was ill I took two days off.
\end{quote}

Feminine.

\mr{रजा} \mr{घेणे}, to take leave; \mr{रजेवर} \mr{असणे}, to be on leave.

\begin{cousinweb}
\textbf{\mr{रजा}} (\emph{rajā}) is a loan from Persian \emph{razā}, \textquotedblleft{}leave, permission\textquotedblright{}, usually traced to Arabic \emph{riḍā}, \textquotedblleft{}consent\textquotedblright{}.
\end{cousinweb}

\begin{grammarlens}[title={What you've built}]
One more word for describing what you lived and felt.
\end{grammarlens}

\subsection*{Your turn}
\begin{itemize}
\raggedright
  \item \emph{Say it:} \emph{rajā}
  \item \emph{Say it:} \emph{rajā}, once more
  \item \emph{Say it:} \emph{ājārī aslyāmuḷe mī don divas rajā ghetlī.} — Because I was ill I took two days off.
  \item \emph{Recall:} say the Marathi for junior, then the Marathi for to strive, then say \emph{rajā} again
\end{itemize}

\begin{tcolorbox}[breakable,colback=teal!4,colframe=teal!35!black,title={Before you move on}]
What does \mr{रजा} mean? (\textquotedblleft{}Time off work; leave\textquotedblright{}.) Say it once more.
\end{tcolorbox}
\section[\texorpdfstring{\mr{बेकारी} (bekārī)}{बेकारी (bekārī)}]{\mr{बेकारी} (bekārī) — unemployment}

\label{lesson:MR-C510-bekari}



\begin{quote}
Before the new one: say the Marathi for to strive, then the Marathi for time off work.
\end{quote}

\subsection*{You'll want to know: \mr{बेकारी}}
\textbf{\mr{बेकारी}} — \emph{bekārī} — \textquotedblleft{}unemployment\textquotedblright{}.

\begin{quote}
\mr{गावात} \mr{बेकारी} \mr{वाढली}, \mr{म्हणून} \mr{तरुण} \mr{शहरात} \mr{गेले}.
\end{quote}

\begin{quote}
\emph{gāvāt bekārī vāḍhlī, mhaṇūn taruṇ śahrāt gele.}
\end{quote}

\begin{quote}
Unemployment grew in the village, so the young people went to the city.
\end{quote}

Feminine.

\begin{cousinweb}
\textbf{\mr{बेकारी}} (\emph{bekārī}) is a loan from Persian \emph{bēkārī}, \textquotedblleft{}idleness\textquotedblright{}, from \emph{bē-}, \textquotedblleft{}without\textquotedblright{}, and \emph{kār}, \textquotedblleft{}work\textquotedblright{}.
\end{cousinweb}

\begin{grammarlens}[title={What you've built}]
One more word for describing what you lived and felt.
\end{grammarlens}

\subsection*{Your turn}
\begin{itemize}
\raggedright
  \item \emph{Say it:} \emph{bekārī}
  \item \emph{Say it:} \emph{bekārī}, once more
  \item \emph{Say it:} \emph{gāvāt bekārī vāḍhlī, mhaṇūn taruṇ śahrāt gele.} — Unemployment grew in the village, so the young people went to the city.
  \item \emph{Recall:} say the Marathi for to strive, then the Marathi for time off work, then say \emph{bekārī} again
\end{itemize}

\begin{tcolorbox}[breakable,colback=teal!4,colframe=teal!35!black,title={Before you move on}]
What does \mr{बेकारी} mean? (\textquotedblleft{}Unemployment\textquotedblright{}.) Say it once more.
\end{tcolorbox}
\section[\texorpdfstring{\mr{राजीनामा} \mr{देणे} (rājīnāmā deṇe)}{राजीनामा देणे (rājīnāmā deṇe)}]{\mr{राजीनामा} \mr{देणे} (rājīnāmā deṇe) — to resign}

\label{lesson:MR-C510-rajinama-dene}



\begin{quote}
Before the new one: say the Marathi for time off work, then the Marathi for unemployment.
\end{quote}

\subsection*{You'll want to know: \mr{राजीनामा} \mr{देणे}}
\textbf{\mr{राजीनामा} \mr{देणे}} — \emph{rājīnāmā deṇe} — \textquotedblleft{}to resign\textquotedblright{}.

\begin{quote}
\mr{मीराने} \mr{जुन्या} \mr{कंपनीचा} \mr{राजीनामा} \mr{दिला}.
\end{quote}

\begin{quote}
\emph{mīrāne junyā kampanīcā rājīnāmā dilā.}
\end{quote}

\begin{quote}
Meera resigned from her old company.
\end{quote}

Past, with -\mr{ने} (he …): \mr{त्याने} \mr{राजीनामा} \mr{दिला}, \emph{tyāne rājīnāmā dilā}.

The job takes -\mr{चा}: \mr{नोकरीचा} \mr{राजीनामा} \mr{देणे}; \mr{राजीनामा} is the letter.

\begin{cousinweb}
\textbf{\mr{राजीनामा} \mr{देणे}} (\emph{rājīnāmā deṇe}) is \mr{राजीनामा}, a loan from Persian \emph{rāżī-nāma}, \textquotedblleft{}a letter of consent\textquotedblright{}, with \mr{देणे}, \textquotedblleft{}to give\textquotedblright{}.
\end{cousinweb}

\begin{grammarlens}[title={What you've built}]
One more word for describing what you lived and felt.
\end{grammarlens}

\subsection*{Your turn}
\begin{itemize}
\raggedright
  \item \emph{Say it:} \emph{rājīnāmā deṇe}
  \item \emph{Say it:} \emph{rājīnāmā deṇe}, once more
  \item \emph{Say it:} \emph{mīrāne junyā kampanīcā rājīnāmā dilā.} — Meera resigned from her old company.
  \item \emph{Recall:} say the Marathi for time off work, then the Marathi for unemployment, then say \emph{rājīnāmā deṇe} again
\end{itemize}

\begin{tcolorbox}[breakable,colback=teal!4,colframe=teal!35!black,title={Before you move on}]
What does \mr{राजीनामा} \mr{देणे} mean? (\textquotedblleft{}To resign\textquotedblright{}.) Say it once more.
\end{tcolorbox}
\section[\texorpdfstring{\mr{निवृत्त} \mr{होणे} (nivṛtta hoṇe)}{निवृत्त होणे (nivṛtta hoṇe)}]{\mr{निवृत्त} \mr{होणे} (nivṛtta hoṇe) — to retire}

\label{lesson:MR-C510-nivrtta-hone}



\begin{quote}
Before the new one: say the Marathi for unemployment, then the Marathi for to resign.
\end{quote}

\subsection*{You'll want to know: \mr{निवृत्त} \mr{होणे}}
\textbf{\mr{निवृत्त} \mr{होणे}} — \emph{nivṛtta hoṇe} — \textquotedblleft{}to retire\textquotedblright{}.

\begin{quote}
\mr{माझे} \mr{वडील} \mr{गेल्या} \mr{वर्षी} \mr{निवृत्त} \mr{झाले}.
\end{quote}

\begin{quote}
\emph{mājhe vaḍīl gelyā varṣī nivṛtta jhāle.}
\end{quote}

\begin{quote}
My father retired last year.
\end{quote}

Past (he …): \mr{तो} \mr{निवृत्त} \mr{झाला}, \emph{to nivṛtta jhālā}.

For a father, out of respect, the plural: \mr{वडील} \mr{निवृत्त} \mr{झाले}.

\begin{cousinweb}
\textbf{\mr{निवृत्त} \mr{होणे}} (\emph{nivṛtta hoṇe}) is the tatsama \mr{निवृत्त}, from Sanskrit \emph{nivṛtta}, \textquotedblleft{}turned back, ceased\textquotedblright{} (Monier-Williams), with \mr{होणे}.
\end{cousinweb}

\begin{grammarlens}[title={What you've built}]
One more word for describing what you lived and felt.
\end{grammarlens}

\subsection*{Your turn}
\begin{itemize}
\raggedright
  \item \emph{Say it:} \emph{nivṛtta hoṇe}
  \item \emph{Say it:} \emph{nivṛtta hoṇe}, once more
  \item \emph{Say it:} \emph{mājhe vaḍīl gelyā varṣī nivṛtta jhāle.} — My father retired last year.
  \item \emph{Recall:} say the Marathi for unemployment, then the Marathi for to resign, then say \emph{nivṛtta hoṇe} again
\end{itemize}

\begin{tcolorbox}[breakable,colback=teal!4,colframe=teal!35!black,title={Before you move on}]
What does \mr{निवृत्त} \mr{होणे} mean? (\textquotedblleft{}To retire\textquotedblright{}.) Say it once more.
\end{tcolorbox}
\section[\texorpdfstring{\mr{बडतर्फ} \mr{करणे} (baḍtarph karṇe)}{बडतर्फ करणे (baḍtarph karṇe)}]{\mr{बडतर्फ} \mr{करणे} (baḍtarph karṇe) — to dismiss}

\label{lesson:MR-C510-badtarph-karne}



\begin{quote}
Before the new one: say the Marathi for to resign, then the Marathi for to retire.
\end{quote}

\subsection*{You'll want to know: \mr{बडतर्फ} \mr{करणे}}
\textbf{\mr{बडतर्फ} \mr{करणे}} — \emph{baḍtarph karṇe} — \textquotedblleft{}to dismiss\textquotedblright{}.

\begin{quote}
\mr{चोरीमुळे} \mr{मालकाने} \mr{त्या} \mr{कामगाराला} \mr{बडतर्फ} \mr{केलं}.
\end{quote}

\begin{quote}
\emph{corīmuḷe mālkāne tyā kāmgārālā baḍtarph kelaṁ.}
\end{quote}

\begin{quote}
The owner dismissed that worker because of the theft.
\end{quote}

Past, with -\mr{ने} (he …): \mr{त्याने} \mr{कामगाराला} \mr{बडतर्फ} \mr{केलं}, \emph{tyāne kāmgārālā baḍtarph kelaṁ}.

The person dismissed takes -\mr{ला}: \mr{कामगाराला} \mr{बडतर्फ} \mr{करणे}.

\begin{cousinweb}
\textbf{\mr{बडतर्फ} \mr{करणे}} (\emph{baḍtarph karṇe}) is \mr{बडतर्फ}, a loan from Persian \emph{bar-taraf}, \textquotedblleft{}set aside, removed\textquotedblright{}, with \mr{करणे}.
\end{cousinweb}

\begin{grammarlens}[title={What you've built}]
One more word for describing what you lived and felt.
\end{grammarlens}

\subsection*{Your turn}
\begin{itemize}
\raggedright
  \item \emph{Say it:} \emph{baḍtarph karṇe}
  \item \emph{Say it:} \emph{baḍtarph karṇe}, once more
  \item \emph{Say it:} \emph{corīmuḷe mālkāne tyā kāmgārālā baḍtarph kelaṁ.} — The owner dismissed that worker because of the theft.
  \item \emph{Recall:} say the Marathi for to resign, then the Marathi for to retire, then say \emph{baḍtarph karṇe} again
\end{itemize}

\begin{tcolorbox}[breakable,colback=teal!4,colframe=teal!35!black,title={Before you move on}]
What does \mr{बडतर्फ} \mr{करणे} mean? (\textquotedblleft{}To dismiss\textquotedblright{}.) Say it once more.
\end{tcolorbox}
